\documentclass{situationist-preprint}

\title{Group Agency}
\msid{ACT-10-a}
\msfamily{ACT-10}
\msdiscipline{Action, Agency and Free Will}
\mstype{Principal manuscript}
\mskeywords{group agency; collective intentionality; discursive dilemma; judgment aggregation; supervenience; corporate agents; methodological individualism}

\begin{document}
\maketitle

\begin{abstract}
Courts, firms, parties and governments are said to believe, decide, intend and act. Individualists hold that such talk is shorthand for what members believe and do; collectivists hold that some groups are agents in their own right. We argue that groups can be agents, with beliefs and intentions not reducible to those of their members. An agent is a system that has representational and motivational states and the capacity to process them and act on them in a way that meets minimal standards of rational unity. We show, using the discursive dilemma of Kornhauser and Sager and the impossibility theorem of List and Pettit, that a group whose attitudes are fixed proposition by proposition from its members' attitudes on the same proposition cannot in general be rationally unified, unless it is a dictatorship. A group that achieves rational unity without dictatorship must therefore form some attitudes that are not a function of its members' attitudes on the same proposition, and may hold attitudes that a majority, or even all, of its members reject. Such a group is an agent, and its attitudes supervene on its members' attitudes only holistically. We trace the history of the debate, analyze further cases, and answer five objections.
\end{abstract}

\section{The question}\label{s:question}

We attribute attitudes and actions to groups as a matter of course. The court held that the contract was binding; the company intends to close the plant; the party believes that the reform will fail; the government declared war. The question is whether these attributions are literally true, and if so, whether what they report is anything over and above what the members believe, intend and do.

Three answers are available. \emph{Eliminativism} holds that the attributions are false or at best figures of speech: only individuals have attitudes, and a sentence about what the company intends is either false or a loose way of reporting what certain individuals intend. \emph{Reductionism} grants that the attributions can be true but holds that their truth consists in facts about members' attitudes of the same kind: the company intends to close the plant just in case its directors, or most of them, intend that it be closed, or some suitable aggregate of such facts obtains. \emph{Realism} about group agency holds that some groups are agents, whose attitudes are fixed by the members' attitudes and actions, since there is nothing else for them to be fixed by, but are not reducible to members' attitudes of the same kind about the same matters.

We defend realism. The argument is the one \citet{ListPettit2011} have made central to the debate, and we state it in a form that makes its premises explicit and separates what it establishes from what it does not. The key observation is that a group that is to function as an agent must satisfy constraints of rational unity, and that, given a theorem about judgment aggregation, these constraints cannot in general be satisfied by a group whose attitudes on each matter are a function of its members' attitudes on that matter alone. Agency requires the group to have attitudes of its own.

\section{A short history of the debate}\label{s:history}

The question has a long history in social theory, jurisprudence and philosophy; we recall the proposals that shaped the current debate, with the consideration that told against each.

\paragraph{Summative accounts.} \citet{Quinton1975} held that statements about the mental states of groups are summative: to say that a group believes that $p$ is to say that all or most of its members believe that $p$. The account is reductionist in the sense of Section~\ref{s:question}. It fails in two directions. A group may fail to believe what all its members believe: every member of a hiring committee may privately believe that a candidate is the strongest, while the committee, having agreed to judge by stated criteria, holds that she is not. And a group may believe what few of its members believe, as the court case of Section~\ref{s:dilemma} shows. The first direction was pressed by \citet{Gilbert1989}, the second by the literature on the doctrinal paradox.

\paragraph{Plural subjects.} \citet{Gilbert1989} proposed that a group believes that $p$ when its members have jointly committed to believe that $p$ as a body: each has openly expressed readiness to be party to such a commitment, and each is thereby bound not to contradict $p$ in her capacity as member. The account is non-summative, since a joint commitment can be formed by members none of whom personally believes $p$. Its limit is that joint commitments are formed one at a time and need not cohere. A body that commits jointly to $p$, to $q$, and, on another occasion, to the negation of their conjunction has three joint commitments and no rational unity. The account explains how a group can have an attitude, but not how a group can be an agent whose attitudes constrain one another.

\paragraph{We-intentions and shared intention.} \citet{Searle1990} held that collective action involves we-intentions, attitudes of the form \enquote{we intend that we do $A$} held in individual heads, and irreducible to individual intentions with the content \enquote{I intend that I do my part}. \citet{Bratman2014} gave a reductive account of shared intention: we intend to $J$ when each of us intends that we $J$ by way of meshing subplans, with common knowledge of these intentions and their interdependence. Both accounts concern small-scale joint action, such as walking together or painting a house together, and both locate the relevant attitudes in the participants. They are accounts of acting together, not of a group agent. They do not apply to a firm that holds an intention none of its present employees formed, or that persists in it as its board is replaced.

\paragraph{Corporate internal decision structures.} \citet{French1979} argued that a corporation is a moral person because it has a \emph{corporate internal decision structure}: an organization chart that assigns roles, together with recognition rules that identify certain decisions, made by persons in certain roles in accordance with corporate policy, as decisions of the corporation. French's account correctly locates group agency in a procedure rather than in coincident attitudes. But he gave no argument that the decisions so identified are not reducible to those of the persons who make them. A procedure on which the chief executive decides everything is a corporate internal decision structure; the resulting \enquote{corporate} decisions are the decisions of one person. \citet{Velasquez1983} argued, against French, that corporate acts are always carried out by members and that intentions belong to those who carry them out, so that only members are responsible. The objection is answered only if it can be shown that there are corporate attitudes which no member's attitudes reproduce.

\paragraph{The doctrinal paradox and the discursive dilemma.} \citet{KornhauserSager1986} observed that a multi-member court deciding a case by majority vote on the outcome may reach a different verdict from one deciding by majority vote on each of the doctrinal premises that jointly determine the outcome. \citet{Pettit2001} generalized the observation into the \emph{discursive dilemma}: any group that forms judgments on logically connected propositions by majority vote on each risks inconsistency, and must choose between responsiveness to its members on each proposition and rational unity across propositions. \citet{ListPettit2002} proved that the dilemma is not a peculiarity of majority voting: no aggregation procedure satisfying a short list of attractive conditions guarantees consistent group judgments. \citet{Pettit2003} and \citet{ListPettit2011} drew the metaphysical consequence that we develop here. A group that resolves the dilemma in favor of rational unity, without handing its judgment to a single member, has attitudes that are not a proposition-by-proposition function of its members' attitudes. Such a group, they argued, is an agent with a mind of its own.

\paragraph{Recent work.} The last fifteen years have seen the realist position extended and contested. \citet{Tollefsen2015} defended group agency on interpretivist grounds, as a matter of the success of the intentional stance \citep{Dennett1987} applied to groups. \citet{Ludwig2016} gave a systematic individualist semantics on which sentences ascribing action to groups are true in virtue of facts about members' actions and intentions, with no group agent in the domain. \citet{Collins2019} argued that some groups, those with decision procedures, can bear duties, while unorganized collectives cannot, and \citet{Lackey2021} argued that group belief is not identical to the joint acceptance of a proposition, since a group may accept for strategic reasons what its evidence does not support. Our argument supports the first three of these in their realism and Lackey's point in its separation of belief from mere acceptance, and is aimed at Ludwig's position in Section~\ref{s:objections}.

\section{Framework}\label{s:framework}

\begin{definition}[Agent]\label{d:agent}
An \emph{agent} is a system that (i) has representational states, which depict how things are; (ii) has motivational states, which specify how it is to make things be; and (iii) has the capacity to process these states and to act on the environment so as to satisfy its motivational states when its representational states depict the environment as unfavorable to them, where (iv) the states satisfy minimal standards of rational unity: the representational states are consistent, and the motivational states are jointly realizable by its own lights.
\end{definition}

Definition~\ref{d:agent} follows \citet[ch.~1]{ListPettit2011}. Clauses (i)--(iii) are the familiar functionalist characterization of a belief-desire system; clause (iv) states the rational unity without which the attribution of attitudes loses its predictive and explanatory point \citep{Dennett1987}. A system whose representational states include both $p$ and not-$p$, or which intends incompatible things, cannot be interpreted as acting on its attitudes, since whatever it does conflicts with one of them. Rational unity is required to a degree, not perfectly; individual agents also fall short. What clause (iv) excludes is a system that is inconsistent \emph{systematically}, as a predictable result of the way its states are formed. We call the representational states of an agent \emph{beliefs} or \emph{judgments}, and its motivational states that issue in plans to which it is committed \emph{intentions}, in the sense of ACT-06, which argues that intentions are plans with commitment.

\begin{definition}[Organized group; agenda; profile]\label{d:group}
An \emph{organized group} is a set of $n \geq 2$ members together with a procedure by which certain attitudes, formed or expressed by members in their capacity as members, determine attitudes of the group on the propositions of an \emph{agenda}. The agenda is a set of propositions closed under negation, together with the logical relations among them. A \emph{profile} is an assignment to each member of a set of attitudes on the agenda. An \emph{aggregation function} maps each admissible profile to a set of group attitudes on the agenda.
\end{definition}

\begin{definition}[Propositionwise reducibility]\label{d:red}
A group's attitudes are \emph{propositionwise reducible} to its members' attitudes if, for every proposition $p$ on the agenda, whether the group holds $p$ is a function of the members' attitudes on $p$ alone. They are \emph{summatively reducible} if, in addition, the group holds $p$ just in case all, or a majority, of members hold $p$.
\end{definition}

Summative reducibility is the content of Quinton's account; propositionwise reducibility is the most general form of the reductionist thesis about attitudes of the same kind about the same matters. It allows any function from members' attitudes on $p$ to the group's attitude on $p$: majority, supermajority, unanimity, weighting by seniority, or the attitude of a designated member. It denies only that the group's attitude on $p$ can depend on members' attitudes on \emph{other} propositions. The thesis of this manuscript, that group beliefs and intentions are \emph{not reducible} to those of members, is the denial of propositionwise reducibility for groups that are agents, together with the denial of the special case in which a single member determines everything.

\begin{definition}[Holistic supervenience]\label{d:holistic}
A group's attitudes \emph{supervene holistically} on its members' attitudes if, for every proposition $p$, whether the group holds $p$ is a function of the entire profile of members' attitudes on the agenda.
\end{definition}

The conditions on aggregation that matter are the following.

\begin{principle}[Robustness]\label{p:robust}
A group agent's procedure must yield group attitudes for every profile of consistent and complete member attitudes on the agenda.
\end{principle}

\begin{principle}[Collective rationality]\label{p:rational}
A group agent's attitudes on its agenda must be consistent, and complete enough to guide action: on each proposition the group must act on, it holds the proposition or its negation.
\end{principle}

\begin{principle}[Non-dictatorship]\label{p:nondict}
For the group to be an agent distinct from any of its members, there must be no member whose attitudes on the agenda fix the group's attitudes in every profile.
\end{principle}

Principle~\ref{p:robust} expresses the requirement that a group agent be able to function whatever its members happen to think; a procedure that breaks down whenever members disagree in certain patterns does not make the group an agent, but at most a sometimes-agent whose functioning depends on luck. Principle~\ref{p:rational} is clause (iv) of Definition~\ref{d:agent} applied to groups. Principle~\ref{p:nondict} is not a condition on agency as such. A dictatorial group may be an agent in a trivial sense, but its agency is that of its dictator, and its attitudes are reducible to his. It is a condition on the interesting case, in which the question of irreducibility arises.

\section{The discursive dilemma}\label{s:dilemma}

The standard illustration is a court of three judges deciding a breach of contract case. The defendant is liable ($r$) just in case there was a valid contract ($p$) and the defendant's action breached it ($q$); the judges accept the doctrine that $r$ is equivalent to $p \wedge q$. Their judgments are as follows.

\begin{center}
\begin{tabular}{lccc}
 & $p$ (contract) & $q$ (breach) & $r$ (liable) \\
\hline
Judge 1 & yes & yes & yes \\
Judge 2 & yes & no & no \\
Judge 3 & no & yes & no \\
\hline
Majority & yes & yes & no \\
\end{tabular}
\end{center}

\begin{lemma}\label{l:dilemma}
Propositionwise majority voting on the agenda $\{p, q, r\}$ with $r \leftrightarrow (p \wedge q)$, and their negations, yields an inconsistent set of group judgments at the profile above, although every member's judgments are consistent.
\end{lemma}

\begin{proof}
Each judge's judgments respect the doctrine: Judge 1 accepts both premises and the conclusion; Judges 2 and 3 each reject one premise and the conclusion. A majority (Judges 1 and 2) accepts $p$; a majority (Judges 1 and 3) accepts $q$; a majority (Judges 2 and 3) rejects $r$. The set $\{p, q, \neg r, r \leftrightarrow (p \wedge q)\}$ is inconsistent.
\end{proof}

\begin{proposition}\label{p:anyn}
For every odd $n \geq 3$, there is a profile of consistent member judgments on the agenda of Lemma~\ref{l:dilemma} at which propositionwise majority voting is inconsistent.
\end{proposition}

\begin{proof}
Let $n = 2k+1$ with $k \geq 1$. Divide the members into three blocks: a block $A$ of one member who accepts $p$, $q$ and $r$; a block $B$ of $k$ members who accept $p$ and reject $q$ and $r$; and a block $C$ of $k$ members who reject $p$, accept $q$ and reject $r$. Every member is consistent. The members accepting $p$ are those in $A \cup B$, numbering $k+1$, a majority of $2k+1$; likewise for $q$ and $A \cup C$. The members accepting $r$ are those in $A$, numbering $1$, which is not a majority since $k \geq 1$. The majority judgments are $p$, $q$, $\neg r$, which are inconsistent with the doctrine.
\end{proof}

Proposition~\ref{p:anyn} shows that the dilemma does not depend on the court's small size. It arises for any body voting by majority on connected propositions, and does so at profiles in which the members are evenly divided in an entirely ordinary way. What Lemma~\ref{l:dilemma} shows about majority voting, the impossibility theorem shows about every propositionwise procedure that does not hand the group to a dictator.

\begin{theorem}[\citealp{ListPettit2002}; \citealp{PaulyVanHees2006}; \citealp{DietrichList2007}]\label{t:imposs}
Let the agenda be sufficiently rich in logical connections, as an agenda containing two atomic propositions and their conjunction, with negations, is. Then there is no aggregation function that satisfies Principle~\ref{p:robust}, yields consistent and complete group judgments (Principle~\ref{p:rational}), makes the group's judgment on each proposition depend only on members' judgments on that proposition (Definition~\ref{d:red}), respects unanimity, and is non-dictatorial (Principle~\ref{p:nondict}).
\end{theorem}

We do not reproduce the proof, which is given in the works cited. Its idea can be seen in the court case. If the group's judgment on each proposition depends only on members' judgments on that proposition, and the procedure treats propositions alike, then whichever coalitions are decisive on $p$ are decisive on $q$ and on $r$. Lemma~\ref{l:dilemma} exhibits two decisive coalitions for the premises, whose intersection is a single member, while the complement of that member is decisive against the conclusion. Consistency then forces decisive coalitions to be closed under intersection, and a family of decisive coalitions closed under intersection, on a finite set of members, contains a smallest one; robustness and unanimity force it to be non-empty, and the argument shows it to be a singleton. Its sole member is a dictator. The versions due to \citet{PaulyVanHees2006} and \citet{DietrichList2007} drop the requirement that propositions be treated alike, at the cost of requiring richer agendas, and the conclusion is the same.

\section{The argument}\label{s:argument}

\begin{lemma}\label{l:forced}
Let a group satisfy Principles~\ref{p:robust}--\ref{p:nondict} and respect unanimity on a sufficiently rich agenda. Then its attitudes are not propositionwise reducible to its members' attitudes: for some proposition $p$ and some profile, the group's attitude on $p$ depends on members' attitudes on other propositions.
\end{lemma}

\begin{proof}
If the group's attitudes were propositionwise reducible, its aggregation function would satisfy all the conditions of Theorem~\ref{t:imposs}, which no function does.
\end{proof}

Lemma~\ref{l:forced} concerns any group that satisfies the conditions. We now show that some possible groups do satisfy them and meet Definition~\ref{d:agent}.

\begin{lemma}\label{l:premise}
The court of Section~\ref{s:dilemma}, operating a premise-based procedure (majority voting on $p$ and on $q$, with the judgment on $r$ derived by the doctrine), forms consistent and complete judgments at every profile, and is not a dictatorship. At the profile of Lemma~\ref{l:dilemma} it judges that the defendant is liable, while two of its three members judge that he is not.
\end{lemma}

\begin{proof}
At every profile, majority voting on the two premises yields a judgment on each, since the number of judges is odd; the judgment on $r$ is fixed by the doctrine from these; the result is consistent and complete. No judge is a dictator: at the profile of Lemma~\ref{l:dilemma}, Judge 2 rejects $q$ and the court accepts it, Judge 3 rejects $p$ and the court accepts it, and in the profile in which Judge 1 rejects $p$ and the other two accept it, the court accepts $p$. At the profile of Lemma~\ref{l:dilemma}, the court accepts $p$ and $q$ by majority and derives $r$, while Judges 2 and 3 reject $r$.
\end{proof}

\begin{theorem}\label{t:main}
Groups can be agents with beliefs and intentions not reducible to those of their members.
\end{theorem}

\begin{proof}
Consider a group organized as the court of Lemma~\ref{l:premise}, extended in the way any functioning body is: it has standing purposes, fixed by its charter, which specify how it is to act given its judgments (to issue a verdict for the plaintiff if it judges the defendant liable, for the defendant otherwise); and it has authorized members, its clerks and officers, who act on its behalf by executing the verdicts. (i) The group has representational states, its judgments on the agenda. (ii) It has motivational states, the purposes of its charter, and it forms intentions, committed plans to issue and enforce particular verdicts, by applying those purposes to its judgments. (iii) It processes these states and acts on them through its authorized members. (iv) By Lemma~\ref{l:premise} its judgments are consistent at every profile, and its intentions, derived from consistent judgments by a consistent charter, are jointly realizable. By Definition~\ref{d:agent}, the group is an agent. By Lemma~\ref{l:premise} it is not a dictatorship, and by Lemma~\ref{l:forced}, applied to an agenda containing the premises and their conjunction, its judgments are not propositionwise reducible to its members' judgments. In particular, at the profile of Lemma~\ref{l:dilemma} it believes that the defendant is liable, while a majority of its members believe he is not; so its beliefs are not summatively reducible either. Its intention to issue a verdict for the plaintiff is derived from that belief, and is held while a majority of its members, acting as individuals on their own judgments, would intend the opposite verdict. So its intentions are not reducible to members' intentions of the same kind about the same matter.
\end{proof}

The proof uses a court, but nothing turns on its being a court. Any body that forms judgments on interconnected matters, sets policy by applying standing purposes to them, and acts through authorized members is in the same situation: a firm deciding whether to adopt a product line given judgments about cost and demand, a party committee setting a platform given judgments about the economy and about electoral prospects, a scientific panel issuing a recommendation given judgments about evidence and about risk.

\begin{corollary}\label{c:holism}
The attitudes of a non-dictatorial group agent supervene on its members' attitudes holistically but not propositionwise.
\end{corollary}

\begin{proof}
The premise-based court's judgment on $r$ is a function of the members' judgments on $p$ and $q$, so its attitudes supervene holistically (Definition~\ref{d:holistic}); by Lemma~\ref{l:forced} they do not supervene propositionwise. The same holds for any non-dictatorial group satisfying the principles, since a group's attitudes are fixed by its procedure applied to the profile.
\end{proof}

Corollary~\ref{c:holism} is the precise sense in which group attitudes are \enquote{not reducible}. It is not the claim that they float free of members, which would be incredible: nothing happens in a court that its judges, clerks and officers do not do. It is the claim that no member, and no proposition-by-proposition aggregate of members' attitudes, has the attitude the group has. The group's attitude is determined by the structure of the procedure as it combines members' attitudes across propositions, and it is that structure, rather than any member's mind, that holds the group's attitudes together and makes them rationally unified.

\section{Further cases}\label{s:cases}

\paragraph{The conclusion-based court.} Suppose the court of Section~\ref{s:dilemma} votes only on the conclusion, and gives no reasons. At the profile of Lemma~\ref{l:dilemma} it finds for the defendant. Its single judgment is summatively reducible, and it is consistent, but only because it forms no judgments on the premises. A body that refuses to form attitudes on the premises of its decisions avoids the dilemma at the cost of agency of a substantial kind: it cannot be held to its reasons, cannot apply them to future cases, and cannot be criticized for inconsistency between decisions that turn on the same premise. \citet{Pettit2001} argued that bodies which must give reasons, such as courts bound by precedent and governments bound to publicly justify their policies, are under pressure to collectivize reason in the premise-based way. The more a body functions as an agent across time, the less it can avoid irreducible attitudes.

\paragraph{The committee that intends what most members oppose.} A university appointments committee of three assesses two candidates, $X$ and $Y$, on research, teaching and service, and its rules require that it appoint the candidate who is better on at least two of the three. Each member applies the rule to her own assessments. The first member judges $X$ better on research and teaching; the second judges $X$ better on research only; the third judges $X$ better on teaching only. The first member's application of the rule selects $X$, the other two select $Y$. If the committee proceeds criterion by criterion, as its rules require, it finds $X$ better on research (members one and two) and on teaching (members one and three), and forms the intention to appoint $X$, whom two of its three members, applying the rule to their own assessments, would not appoint. The structure is that of Lemma~\ref{l:dilemma}. The members implement the decision without hypocrisy, since each accepts the procedure. The intention is the committee's.

\paragraph{The dictatorial firm.} A firm whose founder decides every question, and whose employees execute her decisions, has a corporate internal decision structure in French's sense. It satisfies Principles~\ref{p:robust} and \ref{p:rational}, and it can act through its employees. But it violates Principle~\ref{p:nondict}. Its attitudes are reducible to the founder's. Theorem~\ref{t:main} does not count it as an irreducible group agent, and that is the right verdict: talk of what this firm believes is talk of what its founder believes, reported through the firm. Its status changes when it adopts procedures that weigh the judgments of several officers, as most firms of any size do.

\paragraph{The crowd and the market.} A crowd in which each person runs from a fire, and a market in which prices emerge from many transactions, produce collective outcomes that no one intends. Neither is an agent. Neither has representational states of its own, since neither has a procedure that forms judgments on an agenda, and neither satisfies clause (iv) of Definition~\ref{d:agent}, since nothing in its formation constrains its states to be consistent. Collective outcomes, even systematic ones, are not collective agency.

\paragraph{The changing membership.} A court whose judges are replaced one by one retains its precedents. Its new members are bound by its earlier judgments on the premises, so that its later judgments may depend on attitudes of members no longer serving. This is a second source of irreducibility: the group's attitudes at a time are a function not only of the current profile but of its own history. It also shows that the group agent persists through complete change of membership. MET-20-a argues that institutions, once they have come into existence through collective acceptance, persist as long as they are not collectively rejected by the community whose practices contain their rules. Group agents of the kind considered here are institutions in that sense, and their persistence conditions are those of MET-20-a: a court continues while the practice that constitutes it is not rejected, whoever its judges are.

\section{Objections and replies}\label{s:objections}

\begin{objection}[Individualism]
Every action ascribed to a group is performed by some members, and every attitude ascribed to it is realized in members' states. A complete account of what members think and do leaves nothing for the group to do. Sentences about group action are true, but their truth consists in facts about plural individual action \citep{Ludwig2016}.
\end{objection}

\begin{reply}
We grant that group actions are realized by members' actions and group attitudes by members' attitudes and the procedure that combines them. Corollary~\ref{c:holism} says as much. What the objection must show is that this realization amounts to reducibility, and the argument shows that it does not. There is no member's attitude, and no proposition-by-proposition aggregate of members' attitudes, that is the court's belief that the defendant is liable. An individualist semantics can paraphrase the sentence into one about the members' votes on the premises and the doctrine they share, but the paraphrase does not identify an individual attitude with the content of the group's; it describes the procedure that produces the group's attitude. The relation is that between an agent's beliefs and the states of her neurons. That an individual's beliefs are realized in neural states does not make her neurons believers, nor does it make her beliefs reducible to beliefs of anything else.
\end{reply}

\begin{objection}[Acceptance, not belief]
What a court or committee holds is not believed but accepted: adopted as a basis for action for practical purposes, without the truth-directedness of belief \citep{Wray2001}. Groups can accept and decide, but cannot believe.
\end{objection}

\begin{reply}
Some group attitudes are mere acceptances: a negotiating team may adopt a position it has no reason to think true, for strategic reasons. But the attitudes of the premise-based court are formed in response to considerations bearing on whether a contract existed and whether it was breached, and they are revised if new evidence shows the contrary. They are formed by a procedure designed to track the truth on those questions. This is the functional role of belief, and nothing in it requires that the believer be an individual. We accept the distinction, which \citet{Lackey2021} draws carefully, but it divides group attitudes into kinds; it does not show that no group attitude is a belief. The thesis also covers intentions, and that groups can intend is not in question for this objection.
\end{reply}

\begin{objection}[Consciousness]
Beliefs and intentions are states of conscious subjects, or at least of subjects capable of consciousness. Groups are not conscious. So they cannot believe or intend.
\end{objection}

\begin{reply}
Definition~\ref{d:agent} characterizes agency functionally, and the characterization is the one presupposed by the explanatory use of attitude ascriptions in the social sciences and the law. Whether believing requires the capacity for consciousness is a contested thesis about individual minds, and it is one we need not deny for individuals; it is enough that agency in the sense of Definition~\ref{d:agent} does not require it. The group agent's members are conscious, and their deliberation is what the procedure combines, so that the group's processing of reasons is carried out by conscious subjects even though the group's attitudes are not any of theirs. We make no claim that groups have phenomenal states.
\end{reply}

\begin{objection}[Explanatory exclusion]
If the members' attitudes and the procedure fix the group's attitudes, then every effect of the group's attitudes is already explained at the level of members. Group attitudes are causally and explanatorily idle.
\end{objection}

\begin{reply}
Group-level explanations are robust under variation in the realizing facts in a way member-level explanations are not. That the court found the defendant liable because it judged that there was a contract and a breach is an explanation that holds whatever the distribution of votes that produced those judgments, and whoever the judges were. It predicts the court's later decisions, which are bound by the same premises, even after the judges are replaced. An explanation in terms of the particular judges' attitudes would be less general and would not support these predictions. The pattern is the one familiar from the special sciences, where higher-level states that are multiply realized figure in explanations that are not displaced by explanations at the realizing level \citep[ch.~1]{ListPettit2011}.
\end{reply}

\begin{objection}[Fragility]
Theorem~\ref{t:imposs} relies on the condition that the procedure must work for every profile (Principle~\ref{p:robust}). Real groups deliberate until they agree, and their members' attitudes converge. Restrict the domain to profiles that arise after deliberation, and propositionwise majority voting may be consistent, so that the reductive picture is restored.
\end{objection}

\begin{reply}
There are domain restrictions, such as a shared ordering of the members along a single dimension, under which propositionwise majority voting is consistent \citep{ListPettit2011}. But a group that functions as an agent only when its members' views happen to fall into such a pattern is not robust in the sense that agency requires. Courts, firms and governments must reach decisions when deliberation has not produced convergence, and the cases in which it has not are precisely the cases in which the question of how the group's attitudes relate to members' arises. A group that relied on deliberation to deliver a favorable profile would be agential only by luck.
\end{reply}

\section{What the result does not establish}\label{s:limits}

Theorem~\ref{t:main} is a possibility result: some groups can be agents with irreducible attitudes. It does not establish that every group, or every organized group, is an agent. Many groups lack the procedures and the purposes Definition~\ref{d:agent} requires, and many organizations fail clause (iv) systematically.

Nor does it establish that group agents are responsible for what they do. Responsibility requires more than agency: ACT-01-a locates free agency in an action's issuing from a mechanism that is reasons-responsive and the agent's own, and whether a group's procedure satisfies these conditions is a further question. The companion manuscript ACT-10-b addresses the responsibility of groups. The relevance of the result to the allocation of responsibility is nevertheless direct. ACT-16 argues that responsibility for the outputs of artificial systems lies with those who deploy them, and ACT-16-a observes that deployers may be groups. If group agents exist, responsibility can attach to a deploying firm without being distributed without remainder to its members.

The result does not establish that group agents are persons, that they have moral status, or that their interests count independently of those of their members. A group agent is an agent in the sense of Definition~\ref{d:agent}; the capacities that ground moral status, which ETH-07 identifies with sentience, are not among those the definition requires.

Finally, the result concerns the relation between group judgments and individual judgments, not between group preferences and individual preferences. The analogue for preferences is Arrow's theorem, which POL-17 discusses in the theory of voting rules. The two results have a common structure: aggregation that is local, in the sense of considering one issue at a time, cannot in general produce collective attitudes with the rational unity of an individual's. But Theorem~\ref{t:imposs} does not follow from Arrow's theorem, and does not require that members have preferences at all.

\section{Conclusion}\label{s:conclusion}

A group is an agent if it has representational and motivational states, processes them and acts on them, and meets minimal standards of rational unity. A group that meets these standards robustly and without dictatorship cannot form its attitudes proposition by proposition from its members' attitudes on the same proposition, since no aggregation procedure of that kind is consistent across all profiles. Such a group must form some attitudes holistically, and it can then believe and intend what a majority, or all, of its members reject. The premise-based court is an example, and firms, committees and governments that must reason consistently across interconnected questions are in the same situation. Groups can be agents, with beliefs and intentions that supervene on their members' attitudes but are not reducible to them.

\begin{thebibliography}{99}\small
\bibitem[Bratman(2014)]{Bratman2014} Bratman, M.~E. (2014). \emph{Shared Agency: A Planning Theory of Acting Together}. Oxford University Press.
\bibitem[Collins(2019)]{Collins2019} Collins, S. (2019). \emph{Group Duties: Their Existence and Their Implications for Individuals}. Oxford University Press.
\bibitem[Dennett(1987)]{Dennett1987} Dennett, D.~C. (1987). \emph{The Intentional Stance}. MIT Press.
\bibitem[Dietrich and List(2007)]{DietrichList2007} Dietrich, F. and List, C. (2007). Arrow's theorem in judgment aggregation. \emph{Social Choice and Welfare}, 29(1), 19--33.
\bibitem[French(1979)]{French1979} French, P.~A. (1979). The corporation as a moral person. \emph{American Philosophical Quarterly}, 16(3), 207--215.
\bibitem[Gilbert(1989)]{Gilbert1989} Gilbert, M. (1989). \emph{On Social Facts}. Routledge.
\bibitem[Kornhauser and Sager(1986)]{KornhauserSager1986} Kornhauser, L.~A. and Sager, L.~G. (1986). Unpacking the court. \emph{Yale Law Journal}, 96(1), 82--117.
\bibitem[Lackey(2021)]{Lackey2021} Lackey, J. (2021). \emph{The Epistemology of Groups}. Oxford University Press.
\bibitem[List and Pettit(2002)]{ListPettit2002} List, C. and Pettit, P. (2002). Aggregating sets of judgments: an impossibility result. \emph{Economics and Philosophy}, 18(1), 89--110.
\bibitem[List and Pettit(2011)]{ListPettit2011} List, C. and Pettit, P. (2011). \emph{Group Agency: The Possibility, Design, and Status of Corporate Agents}. Oxford University Press.
\bibitem[Ludwig(2016)]{Ludwig2016} Ludwig, K. (2016). \emph{From Individual to Plural Agency: Collective Action, Volume 1}. Oxford University Press.
\bibitem[Pauly and van Hees(2006)]{PaulyVanHees2006} Pauly, M. and van Hees, M. (2006). Logical constraints on judgement aggregation. \emph{Journal of Philosophical Logic}, 35(6), 569--585.
\bibitem[Pettit(2001)]{Pettit2001} Pettit, P. (2001). Deliberative democracy and the discursive dilemma. \emph{Philosophical Issues}, 11, 268--299.
\bibitem[Pettit(2003)]{Pettit2003} Pettit, P. (2003). Groups with minds of their own. In F.~F. Schmitt (ed.), \emph{Socializing Metaphysics: The Nature of Social Reality}. Rowman \& Littlefield.
\bibitem[Quinton(1975)]{Quinton1975} Quinton, A. (1975). Social objects. \emph{Proceedings of the Aristotelian Society}, 76, 1--27.
\bibitem[Searle(1990)]{Searle1990} Searle, J.~R. (1990). Collective intentions and actions. In P.~R. Cohen, J. Morgan and M.~E. Pollack (eds.), \emph{Intentions in Communication}. MIT Press.
\bibitem[Tollefsen(2015)]{Tollefsen2015} Tollefsen, D.~P. (2015). \emph{Groups as Agents}. Polity.
\bibitem[Velasquez(1983)]{Velasquez1983} Velasquez, M.~G. (1983). Why corporations are not morally responsible for anything they do. \emph{Business and Professional Ethics Journal}, 2(3).
\bibitem[Wray(2001)]{Wray2001} Wray, K.~B. (2001). Collective belief and acceptance. \emph{Synthese}, 129(3), 319--333.
\end{thebibliography}

\end{document}
